\section{Lean formalization architecture}
\label{sec:lean}

\subsection{Target API}

The implementation should extend the checked repository workspace rather than
copying a second logical kernel.  The primary imports are:

\begin{lstlisting}
import ConciseIPRA.PR
import ConciseIPRA.NativeBiInterpretation
\end{lstlisting}

The relevant checked interfaces are:
\begin{itemize}
\item \code{Syn}, \code{Fm}, scoped substitution, shifting, and instantiation;
\item \code{Prf} natural deduction, \code{Prf.subst}, weakening, cut, and
      equality substitution;
\item \code{PR}, \code{PTm}, \code{PRAAx}, \code{PRA.precZero},
      \code{PRA.precSucc}, and \code{PRA.qfInd}; and
\item \code{pra\_iff\_native} for the final transport.
\end{itemize}

Two existing packet modules should be promoted into a common
\code{ConciseIPRA/Finitary} namespace:
\begin{enumerate}[label=(\roman*)]
\item the checked \code{primitiveRecursionInvariant} and
      \code{primitiveRecursionAt} transformer from
      \code{ipra\_local\_global}; and
\item the checked certificate/stable-observable combinators from
      \code{ipra\_stable\_observables}.
\end{enumerate}
Their conclusions already have the required \code{Prf PRAAx} type.

\subsection{Canonical data modules}

The first implementation milestone is one shared representation library:

\begin{center}
\begin{tabularx}{0.94\textwidth}{@{}lX@{}}
\toprule
Module & Contents \\
\midrule
\code{Finitary/Bool.lean} & normalized bits and truth-table proofs \\
\code{Finitary/Pair.lean} & pairing, projections, bounded tuples \\
\code{Finitary/List.lean} & bounded lists, lookup, update, folds \\
\code{Finitary/Int.lean} & signed integers and order \\
\code{Finitary/Rat.lean} & positive-denominator rationals and dyadics \\
\code{Finitary/Table.lean} & bounded maps, bit vectors, path and trace codes \\
\bottomrule
\end{tabularx}
\end{center}

Every operation should have a \(\PR\) symbol, a term application helper, and
embedded defining/correctness derivations.  The same code must be reused across
all application packets.

\subsection{Explicit trace formulas}

The accompanying file \code{TraceReflectionRoadmap.lean} defines the formula
constructors in full.  The intended core is:

\begin{lstlisting}
structure PRTraceSystem where
  good   : PR 1
  init   : PR 1
  bound  : PR 1
  step   : PR 3       -- (time, previous state, input)
  inv    : PR 3       -- (time, state, input)
  accept : PR 2       -- (state, input)

def PRTraceSystem.run (T : PRTraceSystem) : PR 2 :=
  .prec T.init T.step

structure TraceProofs (T : PRTraceSystem) where
  base  : Prf PRAAx [] (BaseFm T)
  step  : Prf PRAAx [] (StepFm T)
  final : Prf PRAAx [] (FinalFm T)
\end{lstlisting}

The main implementation should be a definition computing a derivation:

\begin{lstlisting}
def traceReflection (T : PRTraceSystem) (h : TraceProofs T) :
    Prf PRAAx [] (ResultFm T) := by
  -- Put Good(u) in the context.
  -- Derive the base and step premises for InvP(state,time,u).
  -- Call primitiveRecursionAt at bound(u).
  -- Instantiate h.final at the terminal run state.
  -- Discharge Good(u), then introduce u.
  ...
\end{lstlisting}

The key engineering choice is to call \code{primitiveRecursionAt}; this avoids
repeating its already checked substitution lemmas.  The remaining work is a
small number of \code{simpa} statements identifying the trace formulas with
\code{baseFormula}, \code{nextFormula}, and \code{instStateTime}.

\subsection{Demand compilation}

The object demand module should define fixed symbols for:
\begin{align*}
 \mathsf{backward}(D,N,k,j),\qquad
 \mathsf{horizon}(D,N,k,t),\qquad
 \mathsf{advance}(F,H,s_0,t).
\end{align*}
Because there is no universal evaluator for a variable \(D\), the practical
Lean interface is parameterized by the actual \code{PR} symbols:

\begin{lstlisting}
structure PRDemandSystem where
  demand : PR 2       -- (time, requested output precision)
  base   : PR 1       -- initial object at a precision
  step   : PR 3       -- (time, next precision, previous object)

def backwardSchedule (D : PRDemandSystem) : PR 3 := ...
def compiledRun      (D : PRDemandSystem) : PR 2 := ...
\end{lstlisting}

The schedule equations are proved with \code{PRA.precZero} and
\code{PRA.precSucc}.  The demand-passing theorem then calls the ordinary
primitive-recursion invariant transformer with the formula
\(P(t,\mathsf{horizon}(N,k,t))\).  No new induction rule is introduced.

For variable finite predecessor lists, implement an envelope symbol with nested
finite maxima.  Prove query coverage once; later applications need only supply
the local Booleanized step theorem.

\subsection{Relational kernel}

A numerical layered presentation can be represented by:

\begin{lstlisting}
structure PRLayeredRelation where
  bound  : PR 2       -- (level, parameter)
  good   : PR 3       -- (candidate, level, parameter)
  compat : PR 4       -- (next, candidate, level, parameter)

def viable     (R : PRLayeredRelation) : PR 4 := ...
def reach       (R : PRLayeredRelation) : PR 3 := ...
def fullViable  (R : PRLayeredRelation) : PR 4 := ...
def leastPath   (R : PRLayeredRelation) : PR 5 := ...
\end{lstlisting}

Suggested object endpoints are:
\begin{lstlisting}
def viableIffPath (R : PRLayeredRelation) :
  Prf PRAAx [] (ViablePathSpec R) := ...

def leastViablePathSound (R : PRLayeredRelation) :
  Prf PRAAx [] (LeastPathSpec R) := ...

def fullViableSplit (R : PRLayeredRelation) :
  Prf PRAAx [] (FullViableSpec R) := ...
\end{lstlisting}

Implement a depth-recursive definition first because its proof is shortest.
Then implement forward/backward tables and prove equivalence.  The table version
has better execution complexity, while the recursive version is the reference
specification.

\subsection{Stable observations and shrinking}

The stable-observable module needs only bounded search and the already checked
cut combinators.  Its assumptions are object derivations of nonemptiness and
pairwise stability.  The main theorem substitutes the two least-certificate
terms into the stability proof.

The shrinking module is larger:

\begin{lstlisting}
structure ShrinkingProofs
    (R : PRLayeredRelation) (obs dist eps delta horizon : PR ...) where
  cloudNonempty : Prf PRAAx [] (CloudNonemptyFm ...)
  horizonMono    : Prf PRAAx [] (HorizonMonoFm ...)
  cloudDiameter  : Prf PRAAx [] (CloudDiameterFm ...)
  edgeBound      : Prf PRAAx [] (EdgeBoundFm ...)
  metric         : MetricProofs dist

def shrinkingViability (h : ShrinkingProofs ...) :
  Prf PRAAx [] (ShrinkingBoundFm ...) := ...
\end{lstlisting}

Its internal dependency chain is: least full-viable candidate, least complete
path, path truncation, finite telescoping, triangle inequality.  Each is a
separate reusable theorem.

\subsection{Packet endpoint records}

Completed application packets should export:

\begin{lstlisting}
structure PacketTarget where
  params      : Nat
  assumptions : List (Fm PTm params)
  conclusion  : Fm PTm params
  derivation  : Prf PRAAx assumptions conclusion

def PacketTarget.toNative (T : PacketTarget) :
    NPrf (T.assumptions.map NFm.ofPRA) (NFm.ofPRA T.conclusion) :=
  (pra_iff_native).mp T.derivation
\end{lstlisting}

An incomplete packet should export a record of missing producer or estimate
obligations, not an axiom, \code{opaque} placeholder, or a semantic surrogate.

\subsection{Dependency-ordered implementation plan}

\begin{enumerate}[label=\textbf{M\arabic*.}]
\item Integrate the checked recursion and observable-glue modules.
\item Complete application helpers and the trace formulas; prove
      \code{traceReflection} via \code{primitiveRecursionAt}.
\item Implement normalized Booleans, pairs, and bounded lists.
\item Prove object-level folds, least search, and finite maximum.
\item Implement demand schedules and parameter-changing transfer.
\item Implement path codes, viability, and least viable paths.
\item Implement signed integers, rationals, finite sums, and error ledgers.
\item Prove stable-observable erasure and shrinking viability.
\item Add regular-name arithmetic, guards, finite nets, and robust branches.
\item Migrate packets one at a time, exporting \code{PacketTarget} values.
\end{enumerate}

The first vertical slice should be a finite sum loop whose exact theorem is
transported to \(\NPrf\).  The second should be the square-root interval example
in \cref{sec:examples}; it exercises finite relations, viability, rational
bounds, and a regular-name modulus.

\subsection{Audit discipline}

For each public endpoint, the repository build should record:
\begin{enumerate}[label=(\arabic*)]
\item the pinned ProveIt commit, Lean version, and Mathlib revision;
\item a scan for \code{sorry}, \code{admit}, and newly introduced axioms;
\item \code{\#print axioms} output;
\item a theorem ledger linking the human statement to the Lean name;
\item executable regression tests, clearly separated from universal proofs; and
\item proof-term size and compile-time measurements for the compiler-generated
      derivation.
\end{enumerate}
